% 冒烟测试：验证 csfstyle.sty 在本机 MiKTeX / xelatex 下能否干净编译，
% 并逐个用到每个新盒子与样式，避免“写了但从没编译过”。
\documentclass[withoutpreface,bwprint]{cumcmthesis}

\usepackage{booktabs,graphicx,amsmath,amssymb,float}
\usepackage{csfstyle}

\IfFontExistsTF{Times New Roman}{\setmainfont{Times New Roman}}{\setmainfont{TeX Gyre Termes}}
% 中文：只设 SimSun + AutoFakeBold（实测这一条即可消除粗体缺字告警）。
% 两个坑：① fontspec 认不出族名 "SimKai"，必须写文件名 simkai.ttf；
%        ② 不要写 ItalicFont=SimKai，会报 "The font SimKai cannot be found"
%           并级联出 nullfont 错误。
\IfFontExistsTF{SimSun}{\setCJKmainfont[AutoFakeBold=2.5]{SimSun}}{}
\IfFontExistsTF{SimHei}{\setCJKsansfont{SimHei}}{}

% 必须在文档类设好中文字体之后再调用
\csfsetupcjkbold

\title{csfstyle 冒烟测试}
\begin{document}
\maketitle

\begin{abstract}
本节验证版式壳：彩色章节标题、六个顶会风格盒子、三线表、图注、算法环境。
若本页能编译且中文无缺字、无字体回退告警，说明 P3 版式壳可用。
\csfkeywords{版式壳\quad 顶会风格\quad 冒烟测试}
\end{abstract}

\section{盒子测试}

\begin{csfcontrib}
\begin{itemize}
  \item 把问题重述为带拥堵外部性的分布式资源分配；
  \item 将机制显式化为可证伪假设并以递进实验验证；
  \item 给出排队论闭合估计与 Lagrange 影子价格解释。
\end{itemize}
\end{csfcontrib}

\begin{csfkey}
动态拥塞感知把总疏散时间从 $426.7\pm5.3$ s 降到 $204.1\pm25.7$ s，
达容量下界 $152.2$ s 的 $74\%$。
\end{csfkey}

\begin{csfasm}
出口服务率恒为 $\mu_e=w_e q$。若出口出现堵塞导致有效宽度下降，
则本文低估疏散时间，偏差方向已知（偏乐观）。
\end{csfasm}

\begin{csfprop}[命题 1（容量下界）]
任意可行疏散满足 $T\ge N/\sum_e\mu_e\approx152.2$ s。
\end{csfprop}

\begin{csftake}
出口流量 Gini 由 $0.322$ 降至 $0.056$，说明收益来自均衡而非单纯加速。
\end{csftake}

\begin{csflimit}
下界假设出口全程满负荷，实际不可达；本文结果仍高于下界约 $34\%$。
\end{csflimit}

\section{表格}

\begin{table}[H]
  \centering
  \caption{方法对比（5 次独立运行，最优加粗）}
  \label{tab:smoke}
  \begin{tabular}{lccc}
    \toprule
    方法 & 总疏散时间/s & 吞吐/(人/s) & Gini \\
    \midrule
    最近出口 & $426.7\pm5.3$ & 0.94 & 0.322 \\
    静态拥塞感知 & $426.7\pm5.3$ & 0.94 & 0.322 \\
    \textbf{动态拥塞感知（本文）} & $\mathbf{204.1\pm25.7}$ & \textbf{1.96} & \textbf{0.056} \\
    \bottomrule
  \end{tabular}
\end{table}

\section{算法环境}

\begin{algorithm}[H]
  \caption{拥塞感知动态重分配}
  \label{alg:smoke}
  \KwIn{位置 $\{p_i\}$，出口宽度 $\{w_e\}$，权重 $\lambda$}
  \KwOut{总疏散时间 $T$}
  \ForEach{时间步 $t$}{
    各出口按 $\mu_e=w_e q$ 服务队列\Comment{分数容量累计取整}\;
    计算代价 $c_i(e)=d_i(e)+\lambda Q_e$\;
    $a_i\gets\arg\min_e c_i(e)$\;
    移动并入队\;
  }
\end{algorithm}

\end{document}
